\section{Introduction}
\label{sec:intro}

Physics-Informed Neural Networks \citep[PINNs;][]{raissi2019} fit a
single neural network to a single PDE by backpropagating the residual
of the PDE itself. Two training-time choices control the result and
are notoriously difficult to set: the relative weights of the loss
terms (residual / boundary / initial / data), and the value of any
physical parameter that enters the operator. Both are typically
\emph{tuned}: weights via grid search or self-adaptive scheduling
\citep{mcclenny2020,bischof2021,wang2024}, parameters via outright
retraining at every value the practitioner cares about. Either way, every
change in the residual costs a fresh training run.

This paper removes that cost. We show that a PINN can be trained
\emph{once} so that the same network covers an entire family of
loss-weight balances or an entire range of a physical coefficient,
with no test-time adaptation, no per-$\lambda$ retuning, and no
solver-generated paired training data --- one residual-loss training
run, then a single forward pass per query. Reference solutions
appear only at evaluation, to compute rel-$L^2$ error.

A separate stream of work in scientific machine learning attacks the
same problem from the operator-learning angle: methods like FNO
\citep{li2020fno} and DeepONet \citep{lu2021deeponet} amortise across
a family of PDE instances by treating the right-hand side or the
initial condition as a function-valued input. They produce a single
network that responds to a continuum of inputs --- but they require a
training distribution of \emph{paired} (input, solution) examples
generated by an external solver; the PDE residual itself plays no
role in training.

This paper sits between those two worlds. We adopt the PINN
residual-loss machinery and turn the parameter into a network input,
so a single PINN-style training run produces a continuum of
solutions indexed by $\lambda$. The construction generalises the
loss-conditional networks of \citet{dosovitskiy2020} --- originally
proposed for image-generation loss-mixture amortisation --- to the
parametric-PDE setting. Concretely:

\begin{quote}\itshape
LC-PINN is a residual-loss training run that produces a continuous
family of PDE solutions over a parameter $\lambda$ --- a physical
scalar coefficient or a loss-weight vector --- with no solver-
generated paired training data, no test-time adaptation, and no per-$\lambda$
retuning.
\end{quote}

\noindent The same architecture serves two distinct uses, the
parametric-coefficient mode ($\lambda$ a physical scalar) and the
loss-weight mode ($\lambda$ a weight vector); both are reported as
separate experiment families below.

\paragraph{Contributions.}
\begin{itemize}\setlength{\itemsep}{0pt}
   \item We formulate the LC-PINN
         (\Cref{eq:lc-network,eq:lc-objective}) and prove a
         $\lambda$-invariance result: at a global optimum of the
         $\lambda$-expected loss, the conditional network is
         $p_\lambda$-a.e.\ a residual-minimiser of the parametric
         PDE indexed by~$\lambda$ (\Cref{prop:lc-invariance}). A
         finite-capacity corollary bounds the per-$\lambda$
         suboptimality by the realisable residual budget plus the
         optimisation gap (\Cref{cor:finite-capacity}).
   \item On 1D parametric Helmholtz ($k \!\in\! [1, 10]$) the
         amortised LC-PINN reaches rel-$L^2 \!\sim\! 9\!\times\!10^{-4}$
         on the grid mean and stays within one order of magnitude of
         its own best across the entire parameter range. Per-$k$
         retrained SA-PINN/ReLoBRaLo baselines vary by three orders of
         magnitude across $k$ --- occasionally beating LC by a factor
         of $\sim\!200$ at one $k$, then catastrophically failing at
         another. LC is the only method whose rel-$L^2$ stays within
         one order of magnitude across the entire $k$-range
         (\Cref{sec:results-helmholtz}).
   \item On 2D parametric Helmholtz on the unit square and 3D
         parametric Helmholtz on the unit cube the same construction
         extends without architectural changes: rel-$L^2 \!\sim\!
         2.4\!\times\!10^{-2}$ in 2D ($\sim\!3\times$ better than
         per-$k$ SA-PINN baselines, \Cref{sec:results-helmholtz-2d})
         and $\sim\!1.7\!\times\!10^{-2}$ in 3D ($\sim\!4\times$
         better than per-$k$ SA-PINN, \Cref{sec:results-helmholtz-3d}).
   \item On a 1D stationary Schrödinger problem with a parametric
         harmonic potential ($\alpha \!\in\! [0.5, 10]$), where the
         parameter enters multiplicatively against $u$ rather than as
         a constant zeroth-order coefficient, the same FiLM+L-BFGS
         construction transfers without architectural changes
         (\Cref{sec:results-schrodinger}); this rules out the
         suspicion that the lift is specific to the wave operator.
   \item On viscous Burgers and capillarity-corrected
         Buckley--Leverett, LC-PINN is within a factor of two of
         the strongest single-shot baseline while parameterising
         the entire loss-weight family from one trained network,
         breaking even on training compute after
         $\sim\!4$~retrainings and amortising $\sim\!29\times$ at
         $K\!=\!100$ (\Cref{sec:results-burgers-bl}).
   \item Ablations on the $K$-shot inference budget, per-step
         $\lambda$-sample count, and uniform vs.\ log-uniform
         sampling distributions (\Cref{sec:results-ablations}).
\end{itemize}
